
\subsection*{\SGAInlineRefTarget{sga3:i:subsection:1-6}1.6. Base change in a functor}

For every \(S\in\operatorname{Ob}(\mathscr C)\), there is a canonical functor
\[
  i_S:\mathscr C_{/S}\longrightarrow\mathscr C
\]
defined by \(i_S(f)=T\) when \(f\) is the arrow \(T\to S\). If
\(f:T\to S\) is an object of \(\mathscr C_{/S}\), we denote by
\(i_{T/S}=i_f\) the functor
\[
  i_{T/S}:(\mathscr C_{/S})_{/T}\longrightarrow\mathscr C_{/S},
\]
and there is a commutative diagram:

\begin{center}
  \begin{tikzcd}[column sep=4.7em,row sep=3.2em]
  (\mathscr C_{/S})_{/T}
    \arrow[r,"i_{T/S}"]
    \arrow[dr,"\simeq"']
  &
  \mathscr C_{/S}
    \arrow[r,"i_S"]
  &
  \mathscr C
  \\
  &
  \mathscr C_{/T}
    \arrow[ur,"i_T"']
  &
  {} ,
\end{tikzcd}

\end{center}

Thus, identifying \((\mathscr C_{/S})_{/T}\) with \(\mathscr C_{/T}\), as
we shall henceforth do,
\[
  i_S\circ i_{T/S}=i_T.
\]
Likewise, if \(\mathscr C\) is identified with \(\mathscr C_{/e}\) when
\(\mathscr C\) has a final object \(e\), then
\(i_{S/e}:\mathscr C_{/S}\to\mathscr C_{/e}\) is identified with \(i_S\).

For \(X\in\operatorname{Ob}(\mathscr C)\) (respectively,
\(Y\in\operatorname{Ob}(\mathscr C_{/S})\)), let \(p_S(X)\)
(respectively, \(p_{T/S}(Y)\)) be the object of \(\mathscr C_{/S}\)
(respectively, of \(\mathscr C_{/T}\)), when it exists, defined by
\(X\times S\) (respectively, \(Y\times_S T\)) equipped with its second
projection:

\begin{center}
  \begin{tikzcd}[column sep=6.5em,row sep=3.2em]
  X\times S
    \arrow[d,"p_S(X)"']
  &
  \textup{resp.}
  &
  Y\times_S T
    \arrow[d,"p_{T/S}(Y)"]
  \\
  S
  &&
  T
  \qquad .
\end{tikzcd}

\end{center}

The partially defined functor \(p_S\) (respectively, \(p_{T/S}\)) is called
the \emph{base-change functor}. By the definition of a product
(respectively, a fiber product), it is the adjoint functor of \(i_S\)
(respectively, of \(i_{T/S}\)).\footnote{Translator's note: the Gille--Polo re-edition (2011), Exposé~I, says ``right
adjoint'' and writes out the adjunction. In our formulation: for an object
\(U\to S\) of \(\mathscr C_{/S}\), the morphisms \(U\to X\times S\) over
\(S\) correspond to the morphisms \(U\to X\), so that
\(\operatorname{Hom}_{\mathscr C_{/S}}(U,p_S(X))\simeq
\operatorname{Hom}_{\mathscr C}(i_S(U),X)\); likewise
\(\operatorname{Hom}_{\mathscr C_{/T}}(V,p_{T/S}(Y))\simeq
\operatorname{Hom}_{\mathscr C_{/S}}(i_{T/S}(V),Y)\) for an object \(V\to T\)
of \(\mathscr C_{/T}\), where \(i_{T/S}(V)\) is \(V\) regarded over \(S\)
through \(T\to S\).}
We also write
\[
  p_S(X)=X_S,
  \qquad
  p_{T/S}(Y)=Y_T.
\]

The functor \(i_S\) defines a restriction functor
\[
  i_S^*:\widehat{\mathscr C}\longrightarrow
    \widehat{\mathscr C_{/S}}.
\]
We write
\[
  \mathbf F_S=i_S^*(\mathbf F)=\mathbf F\circ i_S.
\]
Plainly,
\[
  i_{T/S}^*\circ i_S^*=i_T^*,
\]
that is, for every functor
\(\mathbf F\in\operatorname{Ob}(\widehat{\mathscr C})\),
\[
  (\mathbf F_S)_T=\mathbf F_T.
\]
The notation is justified by the following result.

\medskip
\noindent\SGAInlineRefTarget{sga3:i:proposition:1-6-1}\textbf{Proposition 1.6.1.}
\textit{The functor
\[
  (\mathbf h_X)_S:(\mathscr C_{/S})^{\circ}\longrightarrow\mathbf{Set}
\]
is representable if and only if the product \(X\times S\) exists. In that
case
\[
  (\mathbf h_X)_S\simeq\mathbf h_{X_S}.
\]}

This shows that \(\mathbf F_S\) has two interpretations: it is the
restriction of \(\mathbf F\) to \(\mathscr C_{/S}\), and it is the functor
obtained by the base change \(e\leftarrow S\). This leads to the following
notation:

\begin{center}
  \begin{tikzcd}[column sep=3.8em,row sep=3.0em]
  \mathbf F
    \arrow[r,no head]
    \arrow[d,no head]
  &
  \mathbf F_S
    \arrow[r,no head]
    \arrow[d,no head]
  &
  \mathbf F_T
    \arrow[d,no head]
  \\
  e
  &
  S
    \arrow[l]
  &
  T
    \arrow[l]
\end{tikzcd}

\end{center}

This diagram reflects both of the preceding interpretations. Notice that
\[
  \Gamma(\mathbf F_S)
    \simeq\operatorname{Hom}(\mathbf h_S,\mathbf F)
    \simeq\mathbf F(S),
\]
and, in particular,
\[
  \Gamma(X_S)\simeq\operatorname{Hom}(S,X).
\]

\SGARefTarget{sga3:i:subsection:1:7}
\subsection*{1.7. The objects \(\underline{\operatorname{Hom}}\),
\(\underline{\operatorname{Isom}}\), etc.}

Let \(\mathbf F\) and \(\mathbf G\) be two objects of
\(\widehat{\mathscr C}\). We define another object of
\(\widehat{\mathscr C}\) by
\begin{align*}
  \underline{\operatorname{Hom}}(\mathbf F,\mathbf G)(S)
  &={}\operatorname{Hom}_{\widehat{\mathscr C_{/S}}}
       (\mathbf F_S,\mathbf G_S)\\
  &\simeq\operatorname{Hom}_{\widehat{\mathscr C}_{/\mathbf h_S}}
       (\mathbf F\times\mathbf h_S,\mathbf G\times\mathbf h_S)\\
  &\simeq\operatorname{Hom}_{\widehat{\mathscr C}}
       (\mathbf F\times\mathbf h_S,\mathbf G).
\end{align*}
The object \(\underline{\operatorname{Hom}}(\mathbf F,\mathbf G)\) has
the following properties:
\begin{enumerate}
  \item[(i)]
    \(\underline{\operatorname{Hom}}(\underline{\mathbf e},\mathbf G)
      \simeq\mathbf G\).\footnote{Translator's note: the Gille--Polo re-edition (2011), Exposé~I, adds that, for a
third object \(\mathbf E\),
\(\underline{\operatorname{Hom}}(\mathbf E,\mathbf F\times\mathbf G)\simeq
\underline{\operatorname{Hom}}(\mathbf E,\mathbf F)\times
\underline{\operatorname{Hom}}(\mathbf E,\mathbf G)\): at each \(S\), a
morphism \(\mathbf E_S\to\mathbf F_S\times\mathbf G_S\) is a pair of
morphisms.}
  \item[\SGAInlineRefTarget{sga3:i:subsection:1:7:item:ii}(ii)] The formation of \(\underline{\operatorname{Hom}}\) commutes
    with base extension:
    \[
      \underline{\operatorname{Hom}}(\mathbf F_S,\mathbf G_S)
        \simeq
      \underline{\operatorname{Hom}}(\mathbf F,\mathbf G)_S.
    \]
  \item[(iii)]
    \((\mathbf F,\mathbf G)\mapsto
      \underline{\operatorname{Hom}}(\mathbf F,\mathbf G)\)
    is a bifunctor, contravariant in \(\mathbf F\) and covariant in
    \(\mathbf G\).
\end{enumerate}
These three properties follow immediately from the definitions.

Let \(\mathbf E\) be a third object of \(\widehat{\mathscr C}\). We shall
define a morphism of trifunctors
\[
  \operatorname{Hom}(\mathbf E\times\mathbf F,\mathbf G)
    \longrightarrow
  \operatorname{Hom}
    \bigl(\mathbf E,
      \underline{\operatorname{Hom}}(\mathbf F,\mathbf G)\bigr).
\]
Let \(\phi:\mathbf E\times\mathbf F\to\mathbf G\); we must associate
with it a morphism from \(\mathbf E\) to
\(\underline{\operatorname{Hom}}(\mathbf F,\mathbf G)\). So let
\(S'\to S\) be an arrow of \(\mathscr C\). There are maps
\[
  \mathbf E(S)\times\mathbf F(S')
    \longrightarrow
  \mathbf E(S')\times\mathbf F(S')
    \xrightarrow{\ \phi(S')\ }
  \mathbf G(S').
\]
An element of \(\mathbf E(S)\) therefore defines, for every \(S'\to S\),
a map \(\mathbf F(S')\to\mathbf G(S')\) functorial in \(S'\), hence an
element of \(\underline{\operatorname{Hom}}(\mathbf F,\mathbf G)(S)\). We
leave it to the reader to verify that this indeed constructs a morphism of
trifunctors.

\medskip
\noindent\SGAInlineRefTarget{sga3:i:proposition:1-7-1}\textbf{Proposition 1.7.1.}
\textit{The morphism of functors
\[
  \operatorname{Hom}(\mathbf E\times\mathbf F,\mathbf G)
    \longrightarrow
  \operatorname{Hom}
    \bigl(\mathbf E,
      \underline{\operatorname{Hom}}(\mathbf F,\mathbf G)\bigr)
\]
is an isomorphism.}

We indicate only the principle of a proof. Regard both sides as functors
of \(\mathbf E\). The asserted result is true when \(\mathbf E=\mathbf h_X\),
since it is then precisely the definition of the functor
\(\underline{\operatorname{Hom}}(\mathbf F,\mathbf G)\). Moreover, both
sides, as functors of \(\mathbf E\), take direct limits to inverse limits.
Finally, every object of \(\widehat{\mathscr C}\) is a direct limit of
objects of the form \(\mathbf h_X\).

\par\medskip\noindent
\emph{Addition} (after the Gille--Polo re-edition (2011), Exposé~I, 1.7.0
and 1.7.1; our formulation). (a) The last step of the proof. For
\(\mathbf E\in\operatorname{Ob}(\widehat{\mathscr C})\), let
\(\mathscr C_{/\mathbf E}\) be the category of pairs \((V,\rho)\) with
\(V\in\operatorname{Ob}(\mathscr C)\) and \(\rho\in\mathbf E(V)\), that is,
\(\rho:\mathbf h_V\to\mathbf E\), a morphism \((V,\rho)\to(V',\rho')\) being
a morphism \(v:V\to V'\) with \(\mathbf E(v)(\rho')=\rho\). The morphisms
\(\rho\) are compatible and define
\(\psi:\varinjlim_{(V,\rho)}\mathbf h_V\to\mathbf E\), which is an
isomorphism (compare SGA~4, Exposé~I, 3.4). Indeed, the direct limit is
computed pointwise: its value at \(S\) is the set of classes of triples
\((V,\rho,s)\), \(s:S\to V\), for the equivalence relation generated by
\((V,\rho,s)\sim(V',\rho',v\circ s)\), and \(\psi\) sends the class of
\((V,\rho,s)\) to \(\mathbf E(s)(\rho)\). The morphism
\(s:(S,\mathbf E(s)(\rho))\to(V,\rho)\) shows that \((V,\rho,s)\) is
equivalent to \((S,\mathbf E(s)(\rho),\operatorname{id}_S)\); so every class
contains a triple \((S,x,\operatorname{id}_S)\), necessarily with \(x\) the
image of the class, and \(\psi(S)\) is bijective. Consequently
\(\operatorname{Hom}(\mathbf E,\mathbf F)=\varprojlim_{(V,\rho)}\mathbf F(V)\)
for every \(\mathbf F\).

(b) A direct proof of 1.7.1. The morphism of the statement sends \(\phi\)
to the morphism \(x\mapsto\bigl(y\mapsto\phi(S')(x_{S'},y)\bigr)_{S'\to S}\),
where \(x_{S'}\) is the image of \(x\in\mathbf E(S)\) in
\(\mathbf E(S')\). Its inverse sends
\(\theta:\mathbf E\to\underline{\operatorname{Hom}}(\mathbf F,\mathbf G)\)
to \(\phi_\theta\), where \(\phi_\theta(S)(x,y)=\theta(x)_S(y)\) for
\(x\in\mathbf E(S)\) and \(y\in\mathbf F(S)\), and
\(\theta(x)_S:\mathbf F(S)\to\mathbf G(S)\) is the component of
\(\theta(x)\in\underline{\operatorname{Hom}}(\mathbf F,\mathbf G)(S)\) at the
object \(\operatorname{id}_S\) of \(\mathscr C_{/S}\). Starting from
\(\phi\), one recovers \(\phi\) at once; starting from \(\theta\), one
recovers \(\theta\) because, \(\theta\) being a morphism of functors, the
component of \(\theta(x)\) at \(S'\to S\) is \(\theta(x_{S'})_{S'}\).

(c) An internal form. By 1.7 (ii) and 1.7.1 applied to \(\mathscr C_{/S}\),
\[
  \begin{aligned}
  \underline{\operatorname{Hom}}\bigl(\mathbf E,
    \underline{\operatorname{Hom}}(\mathbf F,\mathbf G)\bigr)(S)
  &\simeq\operatorname{Hom}\bigl(\mathbf E_S,
    \underline{\operatorname{Hom}}(\mathbf F_S,\mathbf G_S)\bigr)\\
  &\simeq\operatorname{Hom}(\mathbf E_S\times_S\mathbf F_S,\mathbf G_S)
  =\underline{\operatorname{Hom}}(\mathbf E\times\mathbf F,\mathbf G)(S),
  \end{aligned}
\]
functorially in \(S\). Hence
\(\underline{\operatorname{Hom}}(\mathbf E,\underline{\operatorname{Hom}}(\mathbf F,\mathbf G))
\simeq\underline{\operatorname{Hom}}(\mathbf E\times\mathbf F,\mathbf G)\)
and, exchanging \(\mathbf E\) and \(\mathbf F\),
\(\underline{\operatorname{Hom}}(\mathbf E,\underline{\operatorname{Hom}}(\mathbf F,\mathbf G))
\simeq\underline{\operatorname{Hom}}(\mathbf F,\underline{\operatorname{Hom}}(\mathbf E,\mathbf G))\);
the re-edition refers to this form in Exposé~II, 1 and 3.11.
\par\medskip

\medskip
\noindent\SGAInlineRefTarget{sga3:i:corollary:1-7-2}\textbf{Corollary.}
\textit{One has
\[
  \operatorname{Hom}
    \bigl(\mathbf E,
      \underline{\operatorname{Hom}}(\mathbf F,\mathbf G)\bigr)
    \simeq
  \operatorname{Hom}
    \bigl(\mathbf F,
      \underline{\operatorname{Hom}}(\mathbf E,\mathbf G)\bigr).
\]}

In particular, taking \(\mathbf E=\underline{\mathbf e}\) and using
\(\underline{\operatorname{Hom}}(\underline{\mathbf e},\mathbf G)
\simeq\mathbf G\), we obtain
\[
  \Gamma\bigl(\underline{\operatorname{Hom}}(\mathbf F,\mathbf G)\bigr)
    \simeq\operatorname{Hom}(\mathbf F,\mathbf G).
\]
Composition supplies functorial morphisms
\[
  \underline{\operatorname{Hom}}(\mathbf F,\mathbf G)
    \times
  \underline{\operatorname{Hom}}(\mathbf G,\mathbf H)
    \longrightarrow
  \underline{\operatorname{Hom}}(\mathbf F,\mathbf H).
\]

If \(\mathbf F\) and \(\mathbf G\) are two objects of
\(\widehat{\mathscr C}\), let \(\operatorname{Isom}(\mathbf F,\mathbf G)\)
denote the subset of \(\operatorname{Hom}(\mathbf F,\mathbf G)\) formed by
the isomorphisms from \(\mathbf F\) onto \(\mathbf G\). Define a subobject
\(\underline{\operatorname{Isom}}(\mathbf F,\mathbf G)\) of
\(\underline{\operatorname{Hom}}(\mathbf F,\mathbf G)\) by
\[
  \underline{\operatorname{Isom}}(\mathbf F,\mathbf G)(S)
    =\operatorname{Isom}(\mathbf F_S,\mathbf G_S).
\]
Then
\[
  \Gamma\bigl(\underline{\operatorname{Isom}}
    (\mathbf F,\mathbf G)\bigr)
    \simeq\operatorname{Isom}(\mathbf F,\mathbf G),
  \qquad
  \underline{\operatorname{Isom}}(\mathbf F,\mathbf G)
    \simeq
  \underline{\operatorname{Isom}}(\mathbf G,\mathbf F).
\]
In the special case \(\mathbf F=\mathbf G\), set
\begin{align*}
  \underline{\operatorname{End}}(\mathbf F)
    &=\underline{\operatorname{Hom}}(\mathbf F,\mathbf F),
  &
  \operatorname{End}(\mathbf F)
    &=\operatorname{Hom}(\mathbf F,\mathbf F)
      \simeq\Gamma\bigl(\underline{\operatorname{End}}(\mathbf F)\bigr),\\
  \underline{\operatorname{Aut}}(\mathbf F)
    &=\underline{\operatorname{Isom}}(\mathbf F,\mathbf F),
  &
  \operatorname{Aut}(\mathbf F)
    &=\operatorname{Isom}(\mathbf F,\mathbf F)
      \simeq\Gamma\bigl(\underline{\operatorname{Aut}}(\mathbf F)\bigr).
\end{align*}
Formation of the objects
\(\underline{\operatorname{Hom}}\),
\(\underline{\operatorname{Isom}}\),
\(\underline{\operatorname{Aut}}\), and
\(\underline{\operatorname{End}}\) commutes with base change.

\medskip
\noindent\SGAInlineRefTarget{sga3:i:remark:1-7-3}Note that an object isomorphic to
\(\underline{\operatorname{Isom}}(\mathbf F,\mathbf G)\) may also be
constructed as follows. There is a morphism
\[
  \underline{\operatorname{Hom}}(\mathbf F,\mathbf G)
    \times
  \underline{\operatorname{Hom}}(\mathbf G,\mathbf F)
    \longrightarrow
  \underline{\operatorname{End}}(\mathbf F).
\]
Interchanging \(\mathbf F\) and \(\mathbf G\) gives a morphism
\[
  \underline{\operatorname{Hom}}(\mathbf F,\mathbf G)
    \times
  \underline{\operatorname{Hom}}(\mathbf G,\mathbf F)
    \longrightarrow
  \underline{\operatorname{End}}(\mathbf F)
    \times
  \underline{\operatorname{End}}(\mathbf G).
\]
On the other hand, the identity morphism of \(\mathbf F\) is an element of
\(\operatorname{End}(\mathbf F)\), and hence defines a morphism
\(\underline{\mathbf e}\to
\underline{\operatorname{End}}(\mathbf F)\). Doing the same for
\(\mathbf G\) and taking the product gives a morphism
\[
  \underline{\mathbf e}\longrightarrow
  \underline{\operatorname{End}}(\mathbf F)
    \times
  \underline{\operatorname{End}}(\mathbf G).
\]
It is then immediate that the fiber product of \(\underline{\mathbf e}\)
and
\[
  \underline{\operatorname{Hom}}(\mathbf F,\mathbf G)
    \times
  \underline{\operatorname{Hom}}(\mathbf G,\mathbf F)
\]
over
\(\underline{\operatorname{End}}(\mathbf F)
\times\underline{\operatorname{End}}(\mathbf G)\)
is isomorphic to
\(\underline{\operatorname{Isom}}(\mathbf F,\mathbf G)\).

All these definitions apply, in particular, when
\(\mathbf F=\mathbf h_X\) and \(\mathbf G=\mathbf h_Y\). If
\(\underline{\operatorname{Hom}}(\mathbf h_X,\mathbf h_Y)\) is
representable by an object of \(\mathscr C\), we denote that object by
\(\underline{\operatorname{Hom}}(X,Y)\). It has the following property:
if \(Z\times X\) exists, then
\[
  \operatorname{Hom}
    \bigl(Z,\underline{\operatorname{Hom}}(X,Y)\bigr)
    \simeq
  \operatorname{Hom}(Z\times X,Y).
\]
This property characterizes it when products exist in \(\mathscr C\).

Similarly, one defines, when they exist, the objects
\[
  \underline{\operatorname{Isom}}(X,Y),
  \qquad
  \underline{\operatorname{End}}(X),
  \qquad
  \underline{\operatorname{Aut}}(X).
\]
We merely note that the construction above shows that
\(\underline{\operatorname{Isom}}(X,Y)\) exists whenever fiber products
exist in \(\mathscr C\) and the objects
\[
  \underline{\operatorname{Hom}}(X,Y),
  \quad
  \underline{\operatorname{Hom}}(Y,X),
  \quad
  \underline{\operatorname{End}}(X),
  \quad\text{and}\quad
  \underline{\operatorname{End}}(Y)
\]
exist.

This also applies to categories of the form \(\mathscr C_{/S}\).
The corresponding objects will be denoted as explicitly as possible by
appropriate symbols. For example, if \(T\) and \(T'\) are two objects of
\(\mathscr C\) over \(S\), we denote by
\(\underline{\operatorname{Hom}}_S(T,T')\) the object
\[
  \underline{\operatorname{Hom}}_{\mathscr C_{/S}}(T/S,T'/S),
\]
which will thus be an object of \(\mathscr C\) over \(S\), \ldots

